% Chapter imported from the companion textbook
% (formal-learning-theory-book) with lean annotations added
% to cross-link each formalized statement to its kernel declaration.
% Prose is unchanged; annotations are additive.

% Chapter 5: PAC Learning and the Fundamental Theorem
% fundamental_theorem = Profile B (crown jewel equivalence web),
% agnostic_pac_learning = Profile C (revelation: the ε² price),
% pac_lower_bound = Profile D (negative result),
% nfl_theorem = Profile D (negative result, full proof here),
% realizable/agnostic = Profile A (flags),
% vc_characterization = Profile B (absorbed into fundamental_theorem),
% sauer_shelah_lemma = forward reference to Ch10.

\chapter{PAC Learning and the Fundamental Theorem}
\label{ch:pac}

This is the central chapter of the book.  Every paradigm that follows, online
learning, Gold-style identification, universal learning, will be understood
partly by how it resembles and partly by how it differs from what is
established here.  The chapter builds toward a single destination: the
Fundamental Theorem of Statistical Learning, which characterizes PAC
learnability through a web of nine equivalent conditions.  The path there is
the content.  Each lemma is a foothold; the VC characterization is the summit;
the full equivalence web is the view from the top.

\medskip

The chapter is organized as an ascent.

\begin{enumerate}[leftmargin=*,itemsep=2pt]
\item \textbf{The PAC framework} (\Cref{sec:pac-def}): the definition and its
  moving parts.  Brief, the definition is not the centerpiece.
\item \textbf{The VC characterization proof} (\Cref{sec:vc-char}): the full
  proof that $\VC(\mathcal{H}) < \infty$ if and only if $\mathcal{H}$ is PAC
  learnable, built in four stages through Sauer--Shelah, uniform convergence,
  ERM analysis, and the converse.
\item \textbf{The Fundamental Theorem} (\Cref{sec:fundamental}): all nine
  equivalent conditions, the equivalence web, and the directions we proved versus
  those we cited.
\item \textbf{The agnostic setting} (\Cref{sec:agnostic}): why dropping the
  realizability assumption changes sample complexity from $\Theta(d/\varepsilon)$
  to $\Theta(d/\varepsilon^2)$, and why this gap is not an artifact.
\item \textbf{Lower bounds and No Free Lunch} (\Cref{sec:lower-bounds}): the
  matching lower bound $\Omega(d/\varepsilon)$ and the full NFL proof.
\item \textbf{Computational interlude} (\Cref{sec:computational}): the
  information--computation gap.
\end{enumerate}


% ================================================================
% SECTION 1: THE PAC FRAMEWORK
% ================================================================

\section{The PAC Framework}
\label{sec:pac-def}

Two assumptions frame the discussion.

\begin{defn}[Realizable setting]
\label{def:realizable}
  \lean{BatchLearner}
  \leanok
Learning is \emph{realizable} if the target concept $c^*$ lies in the
hypothesis class: $c^* \in \mathcal{H}$.
\end{defn}

\begin{defn}[Agnostic setting]
\label{def:agnostic}
  \lean{BatchLearner}
  \leanok
Learning is \emph{agnostic} if no assumption is made on the relationship between
$c^*$ and $\mathcal{H}$.  The learner competes with the best hypothesis
$h^* = \arg\min_{h \in \mathcal{H}} \risk_D(h)$.
\end{defn}

The realizable setting is a special case of the agnostic one (take
$h^* = c^*$, achieving zero risk).  The agnostic setting is the one that
matters in practice, but the realizable setting is where the cleanest
characterization lives.  We begin there.

\begin{defn}[PAC Learning {\cite{Valiant1984}}]
\label{def:pac}
  \lean{PACLearnable}
  \leanok
A hypothesis class $\mathcal{H}$ over domain $X$ is \emph{PAC learnable} if
there exists an algorithm $A$ and a function $m_\mathcal{H} \colon
(0,1)^2 \to \N$ such that, for every $\varepsilon, \delta \in (0,1)$ and every
distribution $D$ on $X$:

If $c^* \in \mathcal{H}$ and $S \sim D^m$ with $m \geq
m_\mathcal{H}(\varepsilon, \delta)$, then
\[
  \Pr_{S \sim D^m}\!\bigl[\risk_D(A(S)) > \varepsilon\bigr] \leq \delta.
\]
The function $m_\mathcal{H}(\varepsilon, \delta)$ is the \emph{sample
complexity} of $\mathcal{H}$.  When $m_\mathcal{H}$ is polynomial in
$1/\varepsilon$ and $1/\delta$, the class is \emph{efficiently PAC learnable}
(information-theoretically); computational efficiency is a separate
requirement.
\end{defn}

Before unpacking the definition, we illustrate the full PAC cycle on a
concrete class where every step is visible.

\begin{example}[PAC learning of rectangles: the full cycle]
\label{ex:pac-rectangles}
Let $\mathcal{H}$ be the class of axis-aligned rectangles in $\R^2$:
$h_{(a_1,a_2,b_1,b_2)}(x) = \ind[a_1 \leq x_1 \leq b_1 \text{ and }
a_2 \leq x_2 \leq b_2]$.

\medskip
\noindent\textbf{VC dimension.}  $\VC(\mathcal{H}) = 4$.  Four points,
one near each side of a rectangle, can be shattered: for each subset
$T$ of the four points, the rectangle that tightly encloses exactly
$T$ labels precisely $T$ as positive.  No five points can be shattered:
given any five points in $\R^2$, at least one is ``interior'' (not
extremal in any coordinate), and the labeling that excludes only
that point cannot be realized by a rectangle.

\medskip
\noindent\textbf{ERM and the error region.}  Given sample $S$
consistent with target rectangle $R^*$, the ERM learner returns the
tightest enclosing rectangle $R_S$ of the positive examples.  Since
all positive examples lie in $R^*$, we have $R_S \subseteq R^*$:
zero false positives.  The error region $R^* \setminus R_S$ consists
of four \emph{strips}, one per side of $R^*$, each containing no
positive sample.

\medskip
\noindent\textbf{Sample complexity.}  For each side~$j$ ($j = 1,
\ldots, 4$), let $T_j$ be the strip of $R^*$ closest to side~$j$ with
probability mass exactly $\varepsilon/4$ under~$D$.  If a positive
example falls in $T_j$, then $R_S$ extends into $T_j$ and that
strip's contribution to the error drops below $\varepsilon/4$.  The
probability that strip $T_j$ is missed entirely by $m$ i.i.d.\ samples
is at most $(1 - \varepsilon/4)^m \leq e^{-m\varepsilon/4}$.  A union
bound over four strips gives
\[
  \Pr[\risk_D(R_S) > \varepsilon]
  \;\leq\; 4\, e^{-m\varepsilon/4}.
\]
Setting the right-hand side $\leq \delta$ yields
$m = \frac{4}{\varepsilon}\ln\frac{4}{\delta}
= O\!\bigl(\frac{d + \log(1/\delta)}{\varepsilon}\bigr)$
with $d = 4$, matching the Fundamental Theorem's prediction with the
tight $1/\varepsilon$ dependence.

The analysis makes visible what the general proof obscures: the error
region has \emph{geometric structure} (four strips), the union bound
exploits the \emph{number of strips} (controlled by VC dimension),
and the exponential decay in $m$ comes from the i.i.d.\ assumption
applied to each strip independently.  Replacing $\R^2$ with $\R^k$
gives rectangles with $\VC = 2k$ and $2k$ strips, and the same
argument yields $m = O(k/\varepsilon \cdot \log(k/\delta))$.
\end{example}

Three features of this definition deserve emphasis:

\begin{enumerate}[leftmargin=*,itemsep=3pt]
\item \textbf{Distribution-free.}  The guarantee holds for \emph{every}
  distribution $D$.  The learner does not know $D$ and cannot assume anything
  about it.
\item \textbf{Approximate.}  The learner need not find $c^*$ exactly; error
  $\leq \varepsilon$ suffices.
\item \textbf{Probably.}  The guarantee is probabilistic: it may fail with
  probability $\delta$, but $\delta$ can be driven arbitrarily small by
  drawing more samples.
\end{enumerate}

\begin{remark}[The role of $\delta$]
\label{rmk:delta}
The $\delta$ parameter is structurally unimportant for characterization
purposes: a bound with confidence $1 - \delta$ can always be converted to one
with confidence $1 - \delta'$ by replacing $m$ with $m \cdot
\lceil\log(1/\delta')/\log(1/\delta)\rceil$ and taking a majority vote.
Consequently, the sample complexity depends on $\log(1/\delta)$, not on
$1/\delta$, and the VC characterization is insensitive to how $\delta$ enters.
\end{remark}

\begin{defn}[Empirical Risk Minimization]
\label{def:erm}
  \lean{BatchLearner}
  \leanok
Given sample $S = \{(x_i, c^*(x_i))\}_{i=1}^m$ and hypothesis class
$\mathcal{H}$, the \emph{empirical risk minimizer} is
\[
  \ERM_\mathcal{H}(S) = \arg\min_{h \in \mathcal{H}} \emprisk_S(h),
  \qquad
  \emprisk_S(h) = \frac{1}{m}\sum_{i=1}^m \ind[h(x_i) \neq c^*(x_i)].
\]
In the realizable setting, $\ERM$ returns any $h \in \mathcal{H}$ consistent
with $S$ (since $\emprisk_S(c^*) = 0$).
\end{defn}

The question this chapter answers: \emph{for which $\mathcal{H}$ does ERM
succeed, and when is PAC learning possible at all?}


% ================================================================
% SECTION 2: THE VC CHARACTERIZATION, FULL PROOF
% ================================================================

\section{The VC Characterization}
\label{sec:vc-char}

The main result of this section is:

\begin{thm}[VC Characterization of PAC Learnability {\cite{BlumerEhrenfeuchtHausslerWarmuth1989,VapnikChervonenkis1971}}]
\label{thm:vc-char}
  \lean{fundamental_theorem}
  \leanok
  \uses{def:vc-dim,def:unif-conv,def:pac}
Let $\mathcal{H}$ be a hypothesis class over domain $X$.  The following are
equivalent:
\begin{enumerate}[label=(\roman*)]
\item $\mathcal{H}$ is PAC learnable.
\item $\mathcal{H}$ has the uniform convergence property.
\item $\VC(\mathcal{H}) < \infty$.
\end{enumerate}
Moreover, if $d = \VC(\mathcal{H}) < \infty$, then $\mathcal{H}$ is PAC
learnable with sample complexity
\[
  m_\mathcal{H}(\varepsilon, \delta)
  = O\!\left(\frac{d + \log(1/\delta)}{\varepsilon}\right).
\]
\end{thm}

\begin{remark}[Refinement formalized in this kernel]
\label{rmk:wbvc-refinement}
  \uses{thm:vc-char}
The statement of \Cref{thm:vc-char} as written here is silent on the
measurability of the hypothesis class. The symmetrization argument in
Stage~2 below picks up a two-sided ghost-gap bad event on
$X^m \times X^m$ whose measurability is a separate requirement on
$\mathcal{H}$. The literature proof of
Krapp and Wirth~\cite{KrappWirth2024} states this requirement as a
\emph{Borel} hypothesis on the bad event. This kernel formalizes the
fundamental theorem under a strictly weaker hypothesis: the bad event
need only be \emph{null-measurable} with respect to the completion of
the product measure. The weakening is not cosmetic.
\Cref{chap:measurability} constructs a concept class whose ghost-gap
bad event is analytic but not Borel, and thereby witnesses the strict
gap between the two hypotheses. The Lean identifier
\lean{WellBehavedVCMeasTarget} names the refined hypothesis used in
the kernel statement, and \lean{KrappWirthWellBehaved} names the Krapp
and Wirth version; the implication
\lean{KrappWirthWellBehaved.toWellBehavedVC} sits on one side of the
strict gap, and \Cref{cor:borel-analytic-separation} witnesses the
other.
\end{remark}

The proof proceeds in four stages.  Each stage is a separate lemma,
and each lemma is a foothold on the ascent.

\bigskip
\noindent\textbf{Stage 1:} Finite VC dimension $\Longrightarrow$ polynomial
growth function (Sauer--Shelah).

\noindent\textbf{Stage 2:} Polynomial growth function $\Longrightarrow$
uniform convergence ($\varepsilon$-net/symmetrization argument).

\noindent\textbf{Stage 3:} Uniform convergence $\Longrightarrow$ ERM is a
PAC learner.

\noindent\textbf{Stage 4 (Converse):} Infinite VC dimension
$\Longrightarrow$ not PAC learnable.


%,, , , , , , , , , , , , , , --
% STAGE 1: Sauer--Shelah
%,, , , , , , , , , , , , , , --

\subsection{Stage 1: Sauer--Shelah (Statement)}
\label{sec:sauer-shelah}

The growth function measures the effective richness of $\mathcal{H}$ on finite
samples.

\begin{defn}[Growth Function]
\label{def:growth-function}
For a hypothesis class $\mathcal{H}$ over $X$ and integer $m \geq 1$,
\[
  \growth_\mathcal{H}(m) = \max_{\{x_1, \ldots, x_m\} \subseteq X}
  \bigl|\{(h(x_1), \ldots, h(x_m)) : h \in \mathcal{H}\}\bigr|.
\]
\end{defn}

The growth function satisfies $\growth_\mathcal{H}(m) \leq 2^m$ always, with
equality when $\mathcal{H}$ shatters some set of size $m$.  The Sauer--Shelah
lemma says that finite VC dimension forces a polynomial bound.

\begin{lemma}[Sauer--Shelah {\cite{Sauer1972,Shelah1972}}]
\label{lem:sauer-shelah}
  \lean{sauer_shelah}
  \leanok
  \uses{def:vc-dim,def:growth-function}
If $\VC(\mathcal{H}) = d$, then for all $m \geq d$,
\[
  \growth_\mathcal{H}(m)
  \leq \sum_{i=0}^d \binom{m}{i}
  \leq \left(\frac{em}{d}\right)^d.
\]
\end{lemma}

The proof of the Sauer--Shelah lemma is combinatorial, proceeding by induction
on $m + d$.  It is given in full in the companion textbook's chapter on
Combinatorial Dimensions.\footnote{See \url{https://github.com/Zetetic-Dhruv/formal-learning-theory-book}, Chapter~10.}
For the present chapter, we take it as given
and use only the consequence: if $d = \VC(\mathcal{H}) < \infty$, then
$\growth_\mathcal{H}(m) = O(m^d)$, polynomial, not exponential.


%,, , , , , , , , , , , , , , --
% STAGE 2: Uniform Convergence
%,, , , , , , , , , , , , , , --

\subsection{Stage 2: Uniform Convergence}
\label{sec:uniform-convergence}

The uniform convergence property says that empirical risk converges to true risk
\emph{simultaneously} for all hypotheses in $\mathcal{H}$, not just for a
single fixed $h$.

\begin{defn}[Uniform Convergence]
\label{def:unif-conv}
  \lean{HasUniformConvergence}
  \leanok
A hypothesis class $\mathcal{H}$ has the \emph{uniform convergence property}
if for every $\varepsilon, \delta > 0$, there exists
$m_{\mathrm{UC}}(\varepsilon, \delta)$ such that for all distributions $D$,
whenever $m \geq m_{\mathrm{UC}}(\varepsilon, \delta)$:
\[
  \Pr_{S \sim D^m}\!\left[
    \sup_{h \in \mathcal{H}} |\risk_D(h) - \emprisk_S(h)| > \varepsilon
  \right] \leq \delta.
\]
\end{defn}

\begin{thm}[Uniform Convergence from Finite Growth]
\label{thm:unif-conv}
If $\growth_\mathcal{H}(m) \leq (em/d)^d$ for all $m \geq d$, then
$\mathcal{H}$ has the uniform convergence property with
\[
  m_{\mathrm{UC}}(\varepsilon, \delta)
  = O\!\left(\frac{d \log(1/\varepsilon) + \log(1/\delta)}{\varepsilon^2}\right).
\]
In the realizable setting, the $\varepsilon^2$ denominator improves to
$\varepsilon$, giving $m = O\!\left(\frac{d\log(1/\varepsilon) +
\log(1/\delta)}{\varepsilon}\right)$.
\end{thm}

\begin{proof}
The proof uses the \emph{symmetrization} (or ``ghost sample'') technique of
Vapnik and Chervonenkis~\cite{VapnikChervonenkis1971}.

\medskip
\noindent\textbf{Step 1: Symmetrization.}
Draw two independent samples $S, S' \sim D^m$.  We claim that for
$m \geq 8/\varepsilon^2$,
\[
  \Pr_S\!\left[\sup_{h} |\risk_D(h) - \emprisk_S(h)| > \varepsilon\right]
  \leq 2\,\Pr_{S,S'}\!\left[
    \sup_h |\emprisk_S(h) - \emprisk_{S'}(h)| > \varepsilon/2
  \right].
\]
This follows because if $|\risk_D(h) - \emprisk_S(h)| > \varepsilon$ for some
$h$, then with probability at least $1/2$ (by a Chebyshev argument on $S'$),
the ghost sample satisfies $|\emprisk_{S'}(h) - \risk_D(h)| \leq
\varepsilon/2$, and hence $|\emprisk_S(h) - \emprisk_{S'}(h)| > \varepsilon/2$
by the triangle inequality.

\medskip
\noindent\textbf{Step 2: Permutation argument.}
Let $T = S \cup S'$ be the pooled sample of size $2m$.  Conditioned on $T$,
the partition into $S$ and $S'$ is uniformly random among all
$\binom{2m}{m}$ splits.  By a union bound over the distinct label patterns
that $\mathcal{H}$ induces on $T$:
\[
  \Pr_{S,S'}\!\left[\sup_h |\emprisk_S(h) - \emprisk_{S'}(h)| >
  \varepsilon/2\right]
  \leq \growth_\mathcal{H}(2m) \cdot
  \max_h \Pr_{\text{split}}\!\bigl[|\emprisk_S(h) - \emprisk_{S'}(h)|
  > \varepsilon/2\bigr].
\]
The number of distinct behaviors is at most $\growth_\mathcal{H}(2m)$, and for
each fixed labeling pattern, $\emprisk_S(h) - \emprisk_{S'}(h)$ is a sum of
$2m$ centered random variables (each $\pm 1/m$ depending on which half the
point falls in).

\medskip
\noindent\textbf{Step 3: Hoeffding bound.}
For each fixed label pattern, Hoeffding's inequality gives
\[
  \Pr_{\text{split}}\!\bigl[|\emprisk_S(h) - \emprisk_{S'}(h)| > \varepsilon/2\bigr]
  \leq 2\exp(-m\varepsilon^2/8).
\]
Combining with the Sauer--Shelah bound:
\[
  \Pr_S\!\left[\sup_h |\risk_D(h) - \emprisk_S(h)| > \varepsilon\right]
  \leq 4 \left(\frac{2em}{d}\right)^d \exp\!\left(-\frac{m\varepsilon^2}{8}\right).
\]
Setting this $\leq \delta$ and solving for $m$ gives
\[
  m = O\!\left(\frac{d\log(1/\varepsilon) + \log(1/\delta)}{\varepsilon^2}\right).
\]

\medskip
\noindent\textbf{Realizable improvement.}
In the realizable setting, $\risk_D(c^*) = 0$, so only one-sided deviations
matter: we need $\Pr[\exists h \in \mathcal{H} : \emprisk_S(h) = 0 \text{ but }
\risk_D(h) > \varepsilon] \leq \delta$.  For any fixed $h$ with $\risk_D(h) >
\varepsilon$, the probability that $h$ is consistent with all $m$ samples is at
most $(1 - \varepsilon)^m \leq e^{-m\varepsilon}$.  A union bound over the
$\growth_\mathcal{H}(m)$ distinct behaviors gives
\[
  \Pr[\text{bad}] \leq \left(\frac{em}{d}\right)^d e^{-m\varepsilon}.
\]
Setting this $\leq \delta$ yields $m = O\!\left(\frac{d\log(d/\varepsilon) +
\log(1/\delta)}{\varepsilon}\right)$, confirming the $1/\varepsilon$
dependence.
\end{proof}

\begin{remark}[The Hanneke refinement]
\label{rmk:hanneke}
The optimal sample complexity in the realizable case is
$\Theta\!\left(\frac{d + \log(1/\delta)}{\varepsilon}\right)$, achieved by a
careful one-inclusion graph analysis.  The $\log(d/\varepsilon)$ factor in the
union bound above is a mild overhead; Hanneke~\cite{Hanneke2016} showed it can
be removed entirely, establishing the tight bound.
\end{remark}


%,, , , , , , , , , , , , , , --
% STAGE 3: ERM succeeds
%,, , , , , , , , , , , , , , --

\subsection{Stage 3: Uniform Convergence Implies ERM Succeeds}
\label{sec:erm-succeeds}

\begin{prop}[ERM is a PAC learner under uniform convergence]
\label{prop:erm-pac}
If $\mathcal{H}$ has the uniform convergence property, then $\ERM_\mathcal{H}$
is a PAC learner for $\mathcal{H}$.
\end{prop}

\begin{proof}
Suppose $m \geq m_{\mathrm{UC}}(\varepsilon/2, \delta)$, so that with
probability $\geq 1 - \delta$,
\[
  \sup_{h \in \mathcal{H}} |\risk_D(h) - \emprisk_S(h)| \leq \varepsilon/2.
\]
Let $h_S = \ERM_\mathcal{H}(S)$.  In the realizable case, $\emprisk_S(c^*) =
0$, so $\emprisk_S(h_S) = 0$ as well.  Then:
\[
  \risk_D(h_S)
  = \underbrace{(\risk_D(h_S) - \emprisk_S(h_S))}_{\leq \varepsilon/2}
  + \underbrace{\emprisk_S(h_S)}_{= 0}
  \leq \varepsilon/2 < \varepsilon.
\]
In the general (non-realizable) case, let $h^* = \arg\min_{h \in \mathcal{H}}
\risk_D(h)$.  Then:
\begin{align*}
\risk_D(h_S) &\leq \emprisk_S(h_S) + \varepsilon/2
  \leq \emprisk_S(h^*) + \varepsilon/2
  \leq \risk_D(h^*) + \varepsilon. \qedhere
\end{align*}
\end{proof}


%,, , , , , , , , , , , , , , --
% STAGE 4: Converse
%,, , , , , , , , , , , , , , --

\subsection{Stage 4: Infinite VC Dimension Implies Failure}
\label{sec:converse}

\begin{thm}[Converse of the VC Characterization]
\label{thm:vc-converse}
If $\VC(\mathcal{H}) = \infty$, then $\mathcal{H}$ is not PAC learnable.
\end{thm}

\begin{proof}
The proof constructs, for any candidate learner $A$ and any sample size $m$,
a distribution on which $A$ must fail.  The construction is adversarial.

Since $\VC(\mathcal{H}) = \infty$, there exists a shattered set
$C = \{x_1, \ldots, x_{2m}\} \subseteq X$ of size $2m$.  By the definition of
shattering, for every labeling $b \in \{0,1\}^{2m}$, there exists
$h_b \in \mathcal{H}$ with $h_b(x_i) = b_i$ for all $i$.

Now consider the uniform distribution $D_b$ on $C$ with target concept $h_b$.
Run the learner $A$ on a sample $S$ of size $m$ drawn uniformly from $C$.
With high probability, at least $m$ points of $C$ are unseen
(by a coupon-collector argument, at least $m/2$ are unseen in expectation; we
use $2m$ points to ensure $\geq m$ unseen with constant probability).

On the unseen points, the learner has no information about the target labeling.
We apply a probabilistic argument over $b$ drawn uniformly from $\{0,1\}^{2m}$:

\medskip
\noindent\textbf{Key claim.}  For any fixed algorithm $A$ and fixed training
set $S$, if we draw $b$ uniformly at random, then for any hypothesis $A(S)$
outputs, the expected error on unseen points is at least $1/2 \cdot (m/(2m))
= 1/4$.

This follows because, conditioned on $S$, the labels $b_j$ for unseen points
$x_j \notin S$ are independent uniform bits under the uniform distribution on
target functions.  Therefore $A(S)$'s prediction on each unseen point is
correct with probability exactly $1/2$, regardless of the learner's strategy.

Since the unseen points constitute at least half the support of $D_b$, the
expected true risk satisfies $\E_b[\risk_{D_b}(A(S))] \geq 1/4$.  In
particular, there exists a specific $b^*$ such that
$\risk_{D_{b^*}}(A(S)) \geq 1/4$, which means $A$ fails to achieve
$\varepsilon < 1/4$ with $m$ samples under distribution $D_{b^*}$.  Since $m$
was arbitrary, $\mathcal{H}$ is not PAC learnable.
\end{proof}

\begin{remark}[Strength of the converse]
\label{rmk:converse-strength}
The converse is stronger than ``not uniformly convergent'': it says no learner
whatsoever, not just ERM, can PAC learn $\mathcal{H}$.  The argument works
because the adversary chooses the distribution \emph{after} seeing the learner.
This distribution-free adversarial structure is the engine that makes VC
dimension necessary, not just sufficient.
\end{remark}

The circle closes.  Finite VC dimension forces polynomial growth
(Sauer--Shelah), which forces uniform convergence, which makes ERM
succeed, which gives PAC learnability.  Infinite VC dimension shatters
arbitrarily large sets, and the converse kills learnability outright.
A single combinatorial quantity, the largest set the class can
shatter, controls everything.


% ================================================================
% SECTION 3: THE FUNDAMENTAL THEOREM
% ================================================================

\section{The Fundamental Theorem of Statistical Learning}
\label{sec:fundamental}

The VC characterization (\Cref{thm:vc-char}) is the core equivalence.  But the
full picture is richer: PAC learnability is equivalent to \emph{nine}
conditions, not just three.  The Fundamental Theorem packages them all.

\begin{thm}[The Fundamental Theorem of Statistical Learning
{\cite{BlumerEhrenfeuchtHausslerWarmuth1989,ShalevShwartzBenDavid2014}}]
\label{thm:fundamental}
Let $\mathcal{H}$ be a hypothesis class of binary functions over a domain $X$.
The following are equivalent:
\begin{enumerate}[label=\textnormal{(F\arabic*)}]
\item \label{F:pac} $\mathcal{H}$ is PAC learnable.
\item \label{F:agpac} $\mathcal{H}$ is agnostic PAC learnable.
\item \label{F:erm} $\ERM_\mathcal{H}$ is a PAC learner for $\mathcal{H}$.
\item \label{F:uc} $\mathcal{H}$ has the uniform convergence property.
\item \label{F:finvc} $\VC(\mathcal{H}) < \infty$.
\item \label{F:fingrowth} $\growth_\mathcal{H}(m) < 2^m$ for some $m$.
\item \label{F:compression} $\mathcal{H}$ has a finite compression scheme.
\item \label{F:finiteness} Any finite subclass of $\mathcal{H}$ over a
  finite domain is PAC learnable (with bounds depending on $|\mathcal{H}|$),
  and this property ``lifts'' to $\mathcal{H}$.
\item \label{F:online} $\mathcal{H}$ has finite Littlestone dimension
  $\Longrightarrow$ $\mathcal{H}$ is PAC learnable. \emph{(The converse fails:
  this implication is strict.  See \Cref{ch:separations}.)}
\end{enumerate}
\end{thm}

\begin{remark}[On the numbering]
\label{rmk:numbering}
Condition~\ref{F:online} is an implication, not an equivalence: finite
Littlestone dimension implies finite VC dimension (see the textbook's Chapter~10, Combinatorial Dimensions), but
the converse fails, thresholds on $\R$ have $\VC = 1$ but $\Ldim = \infty$
(\Cref{ex:concept-classes}).  We include it to mark the boundary of the
equivalence web: this is where the PAC--online separation begins.
\end{remark}

\begin{figure}[ht]
\centering
\begin{tikzpicture}[
    node/.style={draw, thick, rounded corners, minimum width=2.4cm,
      minimum height=0.8cm, font=\small, align=center, fill=defblue!8},
    edge/.style={-{Stealth[length=2mm]}, thick},
    bidir/.style={{Stealth[length=2mm]}-{Stealth[length=2mm]}, thick},
    note/.style={font=\tiny, text=notegray, align=center, fill=white,
      inner sep=1pt},
    node distance=1.8cm and 2.2cm
]
% Core nodes
\node[node, fill=thmred!12] (vc) {$\VC(\mathcal{H}) < \infty$\\\ref{F:finvc}};
\node[node, right=2.5cm of vc] (uc) {Uniform\\convergence\\\ref{F:uc}};
\node[node, right=2.5cm of uc] (erm) {ERM is PAC\\\ref{F:erm}};

\node[node, below=2cm of vc] (growth) {$\growth(m) < 2^m$\\\ref{F:fingrowth}};
\node[node, below=2cm of uc, fill=thmred!12] (pac) {PAC\\learnable\\\ref{F:pac}};
\node[node, below=2cm of erm] (agpac) {Agnostic\\PAC\\\ref{F:agpac}};

\node[node, below=2cm of growth] (comp) {Finite\\compression\\\ref{F:compression}};
\node[node, below=2cm of pac] (finite) {Finiteness\\property\\\ref{F:finiteness}};
\node[node, right=2.5cm of finite, draw=edgeorange, thick, fill=edgeorange!8]
  (ldim) {$\Ldim < \infty$\\\ref{F:online}\\(strict)};

% Equivalence edges (bidirectional)
\draw[bidir, defblue] (vc) -- node[note, above] {Sauer--Shelah\\$+$ converse} (uc);
\draw[bidir, defblue] (uc) -- node[note, above] {Prop.~\ref{prop:erm-pac}} (erm);
\draw[bidir, defblue] (vc) -- node[note, left] {Sauer} (growth);
\draw[bidir, defblue] (uc) -- node[note, right] {} (pac);
\draw[bidir, defblue] (pac) -- node[note, above] {} (erm);
\draw[bidir, defblue] (pac) -- node[note, right] {} (agpac);
\draw[bidir, defblue] (growth) -- node[note, left] {} (comp);
\draw[bidir, defblue] (pac) -- node[note, right] {} (finite);

% The strict implication
\draw[edge, edgeorange, dashed] (ldim) -- node[note, below=3pt, text=edgeorange]
  {strict: $\Ldim < \infty \Rightarrow \VC < \infty$\\but not conversely} (finite);

% Main diagonal
\draw[bidir, thmred, very thick] (vc) -- node[note, right, text=thmred]
  {\textbf{Core}} (pac);
\end{tikzpicture}
\caption{The equivalence web of the Fundamental Theorem.  All solid
bidirectional arrows are full equivalences.  The dashed arrow from
$\Ldim < \infty$ is strict: finite Littlestone dimension implies PAC
learnability, but not conversely.  The thick red diagonal marks the VC
characterization, the core equivalence proved in
\Cref{sec:vc-char}.}
\label{fig:equivalence-web}
\end{figure}


\subsection{The Proof Architecture}
\label{sec:proof-architecture}

We have already proved (in \Cref{sec:vc-char}):
\[
  \ref{F:finvc} \Rightarrow \ref{F:uc} \Rightarrow \ref{F:erm}
  \Rightarrow \ref{F:pac}
  \qquad\text{and}\qquad
  \neg\ref{F:finvc} \Rightarrow \neg\ref{F:pac}.
\]
This establishes
$\ref{F:finvc} \Leftrightarrow \ref{F:uc} \Leftrightarrow \ref{F:erm}
\Leftrightarrow \ref{F:pac}$.

The remaining directions:

\begin{itemize}[leftmargin=*,itemsep=4pt]
\item $\ref{F:finvc} \Leftrightarrow \ref{F:fingrowth}$:  By the Sauer--Shelah
  lemma (\Cref{lem:sauer-shelah}), $\VC(\mathcal{H}) = d$ implies
  $\growth_\mathcal{H}(m) \leq (em/d)^d < 2^m$ for $m$ large enough.
  Conversely, if $\VC(\mathcal{H}) = \infty$, then $\growth_\mathcal{H}(m) =
  2^m$ for all $m$ (since $\mathcal{H}$ shatters a set of every finite size).
  This is immediate from the definition of VC dimension.

\item $\ref{F:pac} \Leftrightarrow \ref{F:agpac}$:  Agnostic PAC learnability
  trivially implies realizable PAC learnability (set $c^* \in \mathcal{H}$, so
  the best hypothesis has zero risk).  The converse uses the uniform
  convergence property: if $\mathcal{H}$ has finite VC dimension, then uniform
  convergence holds, and \Cref{prop:erm-pac} shows ERM succeeds in the agnostic
  setting as well (with the $\varepsilon^2$ sample complexity discussed in
  \Cref{sec:agnostic}).

\item $\ref{F:finvc} \Leftrightarrow \ref{F:compression}$:  This is the deepest
  equivalence.  We sketch the argument below.

\item $\ref{F:online} \Rightarrow \ref{F:finvc}$: If $\Ldim(\mathcal{H}) < \infty$,
  then $\VC(\mathcal{H}) \leq \Ldim(\mathcal{H}) < \infty$, since any shattered
  set gives a complete binary mistake tree of the same depth (textbook Chapter~10, Combinatorial Dimensions).
\end{itemize}


\subsection{The Compression Direction (Sketch)}
\label{sec:compression-direction}

\begin{defn}[Compression Scheme]
\label{def:compression}
A \emph{compression scheme} of size $k$ for $\mathcal{H}$ is a pair
$(\kappa, \rho)$ where:
\begin{itemize}[nosep]
\item $\kappa$ (the compressor) maps any sample $S$ consistent with some
  $h \in \mathcal{H}$ to a subsequence $\kappa(S) \subseteq S$ of size
  $\leq k$;
\item $\rho$ (the reconstructor) maps any subsequence of size $\leq k$ to a
  hypothesis $\rho(\kappa(S)) \in Y^X$;
\item $\rho(\kappa(S))$ is consistent with the full sample $S$.
\end{itemize}
\end{defn}

The direction $\ref{F:compression} \Rightarrow \ref{F:pac}$ is classical
and relatively straightforward: a compression scheme of size $k$ yields PAC
learning with sample complexity $O\!\left(\frac{k\log(k/\varepsilon) +
\log(1/\delta)}{\varepsilon}\right)$, since the compressed subsequence encodes
all the information the learner needs, and there are at most $\binom{m}{k}$
possible compressed sets.

The converse direction $\ref{F:finvc} \Rightarrow \ref{F:compression}$ is the
deep one.  A weaker exponential bound was proved by Moran and
Yehudayoff~\cite{MoranYehudayoff2016}; the conjectured linear bound
remains open.

\begin{historical}
\textbf{The sample compression conjecture.}
Littlestone and Warmuth conjectured in 1986 that every class of VC dimension
$d$ has a compression scheme of size at most $d$ (the linear bound).
Moran and Yehudayoff (2016) proved the weaker result that every class of
VC dimension $d$ has a compression scheme of size at most $2^{O(d)}$
(an exponential bound). The original conjecture of a linear bound $O(d)$
remains one of the field's central open problems. This kernel formalizes
the Moran-Yehudayoff exponential bound via an approximate minimax (MWU)
construction, not the conjectured linear bound.
\end{historical}

\begin{proof}[Proof sketch (Moran--Yehudayoff)]
The argument constructs a compression scheme from a
\emph{maximum class} (a class achieving the Sauer--Shelah bound with
equality).  Every class of VC dimension $d$ can be embedded, for sample
complexity purposes, into a maximum class of VC dimension $d$ via a projection
argument.  For maximum classes, the one-inclusion graph has special structure:
its edges can be oriented so that every vertex has in-degree at most $d$.  This
orientation defines a compression scheme, given sample $S$, output the
at-most-$d$ predecessors of the unique vertex in the one-inclusion graph
determined by $S$.  The reconstruction uses the structure of the maximum class
to recover a consistent hypothesis from these $d$ points.

Extending from maximum classes to arbitrary classes of finite VC dimension
requires a more delicate argument involving fractional covers of the
one-inclusion hypergraph, which blows up the size to $2^{O(d)}$.
\end{proof}


% ================================================================
% SECTION 4: THE AGNOSTIC SETTING AND THE ε² PRICE
% ================================================================

\section{The Agnostic Setting and the \texorpdfstring{$\varepsilon^2$}{epsilon-squared} Price}
\label{sec:agnostic}

\begin{defn}[Agnostic PAC Learning]
\label{def:agnostic-pac}
A hypothesis class $\mathcal{H}$ is \emph{agnostic PAC learnable} if there
exists an algorithm $A$ and a function $m_\mathcal{H}(\varepsilon, \delta)$
such that for every $\varepsilon, \delta \in (0,1)$ and every distribution $D$
on $X \times \{0,1\}$ (not necessarily generated by any $c^* \in \mathcal{H}$),
whenever $m \geq m_\mathcal{H}(\varepsilon, \delta)$:
\[
  \Pr_{S \sim D^m}\!\bigl[
    \risk_D(A(S)) - \min_{h \in \mathcal{H}} \risk_D(h) > \varepsilon
  \bigr] \leq \delta.
\]
\end{defn}

The Fundamental Theorem tells us that agnostic PAC learnability is
equivalent to realizable PAC learnability (\ref{F:pac} $\Leftrightarrow$
\ref{F:agpac}).  The same classes are learnable.  But the \emph{sample
complexity} changes, and this change is not cosmetic.

\begin{thm}[The Agnostic Sample Complexity Gap]
\label{thm:agnostic-gap}
Let $d = \VC(\mathcal{H})$.
\begin{enumerate}[label=(\alph*)]
\item In the realizable setting:
  $m(\varepsilon, \delta) = \Theta\!\left(\frac{d +
  \log(1/\delta)}{\varepsilon}\right)$.
\item In the agnostic setting:
  $m(\varepsilon, \delta) = \Theta\!\left(\frac{d +
  \log(1/\delta)}{\varepsilon^2}\right)$.
\end{enumerate}
The gap from $1/\varepsilon$ to $1/\varepsilon^2$ is tight: no algorithm can
achieve $o(d/\varepsilon^2)$ in the agnostic setting, and no algorithm needs
$\omega(d/\varepsilon)$ in the realizable setting.
\end{thm}

Where does the gap come from?  It is tempting to blame the proof technique, perhaps a
cleverer argument could recover $1/\varepsilon$ in the agnostic case.  It
cannot.  The gap is fundamental, and the reason is information-theoretic.

\begin{proof}[Proof sketch of the agnostic lower bound]
Consider the class of threshold functions on $[0,1]$ (VC dimension 1, so
$d = 1$ for simplicity).  In the realizable case, a single misclassified
example pins down the threshold to within $\varepsilon$-precision, requiring
$O(1/\varepsilon)$ samples.

In the agnostic case, the distribution $D$ on $(X, Y)$ may have a base error
rate $\eta = \min_h \risk_D(h) > 0$, and the learner must distinguish between
two hypotheses $h_1, h_2$ whose risks differ by $\varepsilon$: $\risk_D(h_1) =
\eta$ and $\risk_D(h_2) = \eta + \varepsilon$.  Distinguishing these requires
detecting a difference $\varepsilon$ against a background noise level $\eta$.

This is a hypothesis testing problem.  By standard lower bounds (Le Cam's
method or Fano's inequality), distinguishing distributions at total variation
distance $O(\varepsilon)$ from $m$ samples requires
$m = \Omega(1/\varepsilon^2)$.  The noise ``contaminates'' the signal: in the
realizable case, an error is always informative (it eliminates hypotheses),
but in the agnostic case, each error might be noise or signal, and separating
the two costs the extra $1/\varepsilon$ factor.
\end{proof}

\begin{remark}[The revelation]
\label{rmk:revelation}
The $\varepsilon$ vs.\ $\varepsilon^2$ gap is the first structural surprise of
statistical learning theory that cannot be anticipated from finite-class
arguments.  For finite $|\mathcal{H}|$, both realizable and agnostic sample
complexities have the same dependence on $\varepsilon$ (up to the
$\log|\mathcal{H}|$ factor).  The gap emerges only when $\mathcal{H}$ is
infinite and VC dimension replaces $\log|\mathcal{H}|$.  Realizability is
not just a simplifying assumption, it provides a qualitatively different
information structure.
\end{remark}


% ================================================================
% SECTION 5: LOWER BOUNDS AND NO FREE LUNCH
% ================================================================

\section{Lower Bounds and the No-Free-Lunch Theorem}
\label{sec:lower-bounds}

\subsection{The PAC Lower Bound}
\label{sec:pac-lower}

The upper bound of \Cref{thm:vc-char} gives $m = O(d/\varepsilon)$ in the
realizable case.  The following shows this is tight up to constants.

\begin{thm}[PAC Lower Bound]
\label{thm:pac-lower}
For any hypothesis class $\mathcal{H}$ with $\VC(\mathcal{H}) = d \geq 1$, any
PAC learner for $\mathcal{H}$ requires sample complexity
\[
  m(\varepsilon, \delta) = \Omega\!\left(\frac{d}{\varepsilon}\right)
\]
for $\varepsilon \leq 1/8$ and $\delta \leq 1/7$.
\end{thm}

\begin{proof}
Since $\VC(\mathcal{H}) = d$, there exists a shattered set
$C = \{x_1, \ldots, x_d\}$.  For each subset $T \subseteq C$, let
$h_T \in \mathcal{H}$ be the hypothesis that labels exactly $T$ as positive.
This gives $2^d$ distinct hypotheses.

Fix $\varepsilon \leq 1/8$.  For each $T \subseteq C$ with $|T| =
\lfloor d/2 \rfloor$, define the distribution $D_T$: uniform on $C$, with
target $h_T$.  The learner receives $m$ i.i.d.\ samples from $D_T$.

Each sample reveals the label of one point in $C$.  After $m$ samples, the
expected number of unseen points in $C$ is $d(1 - 1/d)^m \geq d \cdot
e^{-2m/d}$ (for $m \leq d$).  For $m \leq d/(8\varepsilon)$, the number of
unseen points is at least $d/4$ in expectation.

On each unseen point, the learner must guess the label.  Since the target is
consistent with both labels for unseen points (by the shattering property),
no strategy can beat chance on unseen points.  The error on each unseen point
contributes $1/d$ to the total risk (since $D_T$ is uniform on $C$).  With
at least $d/4$ unseen points in expectation, the expected risk is at least
$1/4 > \varepsilon$.  A Markov argument converts this to a high-probability
statement, showing $m = \Omega(d/\varepsilon)$.
\end{proof}


\subsection{The No-Free-Lunch Theorem (Full Proof)}
\label{sec:nfl-full}

The companion textbook's Chapter~1 (The Objects of Learning)\footnote{\url{https://github.com/Zetetic-Dhruv/formal-learning-theory-book}, Chapter~1.} gave the statement and a brief argument.
Here we give the full proof with the averaging argument over all target
functions.

\begin{thm}[No Free Lunch, Full Version]
\label{thm:nfl-full}
Let $|X| \geq 2m$.  For any learning algorithm $A$,
\[
  \max_{c \in \{0,1\}^X}\;
  \E_{S \sim D_c^m}\!\bigl[\risk_{D_c}(A(S))\bigr] \geq \frac{1}{4},
\]
where $D_c$ is the uniform distribution on $X$ with target $c$.
\end{thm}

\begin{proof}
Instead of proving the $\max$, we prove the stronger statement with $\E_c$
(expectation over a uniform random target), which implies the $\max$ is at
least as large.

Fix any algorithm $A$.  Let $S = (x_1, \ldots, x_m)$ be drawn uniformly from
$X$ (i.i.d.), and let $c$ be drawn uniformly from $\{0,1\}^X$.  Write
$T = X \setminus \{x_1, \ldots, x_m\}$ for the unseen points.

\medskip
\noindent\textbf{Key observation.}  Conditioned on $S$ and on the labels
$(c(x_1), \ldots, c(x_m))$, the labels $\{c(x) : x \in T\}$ are still
independent uniform bits.  This is because the prior on $c$ is product
measure.

Therefore, for each $x \in T$, regardless of what $A(S)$ predicts:
\[
  \Pr_{c}[A(S)(x) \neq c(x) \mid S, c|_S] = \frac{1}{2}.
\]
The expected risk on unseen points is:
\[
  \E_c\!\left[\frac{1}{|X|}\sum_{x \in T} \ind[A(S)(x) \neq c(x)]
  \;\middle|\; S\right]
  = \frac{|T|}{|X|} \cdot \frac{1}{2}
  \geq \frac{1}{2} \cdot \frac{1}{2} = \frac{1}{4},
\]
using $|T| \geq |X|/2$ (since $m \leq |X|/2$).  Taking expectation over $S$
preserves the bound.
\end{proof}

\begin{remark}[NFL and inductive bias, revisited]
\label{rmk:nfl-revisited}
The NFL theorem is the structural reason that the VC characterization involves
a \emph{restriction} on $\mathcal{H}$.  Without restricting to a class of
finite VC dimension, no learning is possible.  The Fundamental Theorem can
therefore be read as: ``inductive bias (finite VC dimension) is both necessary
and sufficient for statistical learning.''
\end{remark}


% ================================================================
% SECTION 6: COMPUTATIONAL INTERLUDE
% ================================================================

\section{Computational Interlude}
\label{sec:computational}

The Fundamental Theorem is an \emph{information-theoretic} characterization:
it says when enough data exists to learn, regardless of computational cost.
The computational question, when can learning be done in polynomial
time, is a separate and largely open problem.

\begin{thm}[Computational Hardness {\cite{KearnsValiant1994}}]
\label{thm:kearns-valiant}
Under standard cryptographic assumptions (specifically, the existence of
one-way functions), there exist concept classes $\mathcal{C}$ with
$\VC(\mathcal{C}) = O(\log n)$ that are PAC learnable
(information-theoretically) but not \emph{efficiently} PAC learnable (no
polynomial-time algorithm achieves PAC guarantees).
\end{thm}

This result separates the information-theoretic and computational landscapes of
PAC learning.  The Fundamental Theorem characterizes the information boundary;
it says nothing about the computational one.

\begin{computational}
The gap between information and computation is the subject of
the textbook's Chapter~16 (Computational Hardness).\footnote{\url{https://github.com/Zetetic-Dhruv/formal-learning-theory-book}, Chapter~16.}
The key examples:
\begin{itemize}[nosep]
\item \textbf{Properly learning DNF formulas} is computationally hard under
  cryptographic assumptions, even though DNF has polynomial VC dimension.
\item \textbf{Improperly learning DNF} can be done efficiently via boosting
  (using a richer hypothesis class).
\item \textbf{Learning intersections of halfspaces} is hard even improperly,
  under stronger assumptions.
\end{itemize}
The proper/improper distinction (\Cref{def:hypothesis-target}) is
computationally sharp even when it is information-theoretically irrelevant.
\end{computational}


% ================================================================
% SECTION 7: WHAT THIS CHAPTER ESTABLISHED
% ================================================================

\section{What This Chapter Established}
\label{sec:ch5-summary}

The central achievement is the Fundamental Theorem (\Cref{thm:fundamental}),
which says that the following are all the same property of a binary hypothesis
class $\mathcal{H}$:

\begin{center}\small
\renewcommand{\arraystretch}{1.2}
\begin{tabular}{lll}
\toprule
\textbf{Condition} & \textbf{First established} & \textbf{Proof location} \\
\midrule
Finite VC dimension & Vapnik--Chervonenkis (1971) & \Cref{sec:vc-char} \\
Uniform convergence & Vapnik--Chervonenkis (1971) & \Cref{thm:unif-conv} \\
ERM is a PAC learner & Blumer et al.\ (1989) & \Cref{prop:erm-pac} \\
PAC learnable & Valiant (1984) & \Cref{thm:vc-char} \\
Agnostic PAC learnable & Haussler (1992) & \Cref{sec:agnostic} \\
$\growth(m) < 2^m$ for some $m$ & Sauer--Shelah (1972) & \Cref{sec:proof-architecture} \\
Finite compression scheme & Moran--Yehudayoff (2016) & \Cref{sec:compression-direction} \\
\bottomrule
\end{tabular}
\end{center}

Four structural lessons emerge:

\begin{enumerate}[leftmargin=*,itemsep=3pt]
\item \textbf{VC dimension is the right measure.}  Not because it was first, but
  because it sits at the center of a seven-way equivalence.

\item \textbf{The $\varepsilon^2$ price is real.}  The agnostic setting costs
  $1/\varepsilon^2$ where the realizable setting costs $1/\varepsilon$.  This is
  not a proof artifact, it is a lower bound.

\item \textbf{ERM is canonical but not unique.}  The Fundamental Theorem says ERM
  works whenever anything works, but other algorithms (compression-based,
  Bayesian) also achieve PAC guarantees under the same conditions.

\item \textbf{Information $\neq$ computation.}  The Fundamental Theorem
  characterizes when learning is \emph{possible}; the Kearns--Valiant barrier
  shows this does not determine when learning is \emph{efficient}.  The
  computational landscape is the subject of the textbook's Chapter~16 (Computational Hardness).
\end{enumerate}

\begin{historical}
\textbf{Timeline.}  Vapnik and Chervonenkis introduced the VC dimension and
proved the uniform convergence direction in 1971.  Valiant formalized PAC
learning in 1984, without the VC connection.  Blumer, Ehrenfeucht, Haussler,
and Warmuth closed the loop in 1989, proving that finite VC dimension is both
necessary and sufficient for PAC learning.  The compression equivalence was
conjectured by Littlestone and Warmuth (1986) and proved by Moran and
Yehudayoff (2016), a 30-year gap that reflects the depth of the combinatorial
argument.  The tight realizable sample complexity $\Theta(d/\varepsilon)$
(removing logarithmic factors) was established by
Hanneke~\cite{Hanneke2016}.
\end{historical}


% ================================================================
% EXERCISES
% ================================================================

\subsection*{Exercises}

\begin{enumerate}[leftmargin=*,itemsep=6pt]

\item \textbf{VC dimension of convex polygons.}
  Let $\mathcal{H}_k$ be the class of convex $k$-gons in $\R^2$:
  $h_P(x) = 1$ iff $x$ lies inside a convex polygon~$P$ with at
  most $k$ vertices.  Prove that $\VC(\mathcal{H}_k) = 2k + 1$.

  \emph{Lower bound:} Place $2k+1$ points in convex position
  (on a circle).  For any subset $T$ of these points with
  $|T| \leq k$, a convex $k$-gon can be drawn to include exactly
  $T$.  For larger subsets, use the complement labeling and the
  observation that $2k+1 - |T| \leq k$ of the remaining points
  can be avoided.

  \emph{Upper bound:} Show that any $2k+2$ points include a
  labeling unrealizable by a $k$-gon.  Argue that a convex
  polygon with $k$ vertices can ``separate'' at most $2k$
  contiguous arcs on any circle through the points, and $2k+2$
  points in convex position create $2k+2$ arcs.

\item \textbf{The distribution-free assumption is essential.}
  The Fundamental Theorem characterizes \emph{distribution-free}
  PAC learnability.  Show that the equivalence breaks if
  ``distribution-free'' is replaced by ``distribution-dependent.''
  \begin{enumerate}[label=(\alph*)]
  \item Construct a class $\mathcal{H}$ with $\VC(\mathcal{H}) =
    \infty$ and a specific distribution $D$ such that $\mathcal{H}$
    is PAC learnable under $D$ with $m(\varepsilon, \delta) = O(1)$
    samples.
  \item More subtly: construct a class $\mathcal{H}$ with
    $\VC(\mathcal{H}) = \infty$ that is PAC learnable under
    \emph{every} fixed distribution $D$ (with sample complexity
    depending on $D$), yet is not distribution-free PAC learnable.
    \emph{Hint:}  Let $\mathcal{H}$ be all measurable functions on
    $[0,1]$.  For any fixed $D$, the metric entropy of $\mathcal{H}$
    under $L^1(D)$ is finite at every scale, one can $\varepsilon$-net
    the class with $O(1/\varepsilon)$ functions, so PAC learning
    under $D$ is possible.  The distribution-free failure comes from
    the adversary's ability to concentrate $D$ on the points where
    the learner has not yet gathered information.
  \end{enumerate}

\item \textbf{Tight agnostic lower bound via Assouad's lemma.}
  Prove the lower bound
  $m(\varepsilon, \delta) = \Omega(d/\varepsilon^2)$ for agnostic
  PAC learning of any class $\mathcal{H}$ with $\VC(\mathcal{H}) = d$,
  using Assouad's lemma rather than Le~Cam's method.

  \emph{Construction:} Let $C = \{x_1, \ldots, x_d\}$ be a shattered
  set.  For each $b \in \{0,1\}^d$, define $D_b$ as the distribution
  that places mass $1/(2d)$ on each $x_i$ and mass $1/2$ on a
  ``noise point'' $z \notin C$ with label drawn from
  $\mathrm{Bernoulli}(1/2)$.  The target function labels $x_i$ as
  $b_i$ and $z$ as~$1$.  The optimal risk under $D_b$ is $1/4$ (from
  the noise at~$z$).

  Show that any algorithm distinguishing $D_b$ from $D_{b'}$ (where
  $b$ and $b'$ differ in one coordinate) with advantage $\varepsilon$
  requires $\Omega(1/\varepsilon^2)$ samples from the signal at
  a single $x_i$.  Since there are $d$ coordinates, Assouad's
  lemma gives the combined bound $\Omega(d/\varepsilon^2)$.  Verify
  that this matches the agnostic gap of \Cref{thm:agnostic-gap}.

\end{enumerate}

